%! Multiplying \& Dividing Rational Expressions — Unit 4, Lesson 4.2 — Lesson Plan
\documentclass[10pt]{article}
\usepackage{algebra2-article}
\usepackage{algebra2-boxes}
\usepackage{graphicx}   % -article does NOT load graphicx; needed for thumbnails/screenshots
\graphicspath{{images/}}
\newcommand{\TallMath}[1]{$\displaystyle #1\rule[-1.4em]{0pt}{3.2em}$}
\newcommand{\CourseName}{Algebra 2}
\newcommand{\MeetingLength}{60 minutes}
\newcommand{\UnitNumberName}{Unit 4: Rational Functions \quad}
\newcommand{\LessonNumberName}{Lesson 4.2: Multiplying \& Dividing Rational Expressions}

\begin{document}
\begin{center}
    {\Large\bfseries \CourseName \\
    {\small\bfseries \UnitNumberName \LessonNumberName}}
\end{center}

\begin{tcolorbox}[colback=forestbg, colframe=forest, boxrule=0.9pt, arc=2mm,
    left=2mm, right=2mm, top=1.5mm, bottom=1.5mm]
    \textbf{Primary Objective:} Students can multiply and divide rational expressions --- factoring
    every numerator and denominator \emph{first}, dividing out common factors across the whole
    product, and using \textbf{keep--change--flip} (multiplication by the reciprocal) for a quotient,
    flipping \emph{before} cancelling anything. Above all, students can state \emph{every} domain
    restriction from all \textbf{three} sources: each original denominator \emph{and the numerator of
    the divisor}, since a quotient $\frac{A}{B}\div\frac{C}{D}$ requires $B\neq0$, $D\neq0$, and
    $C\neq0$. They explain why source~2 goes invisible \emph{after} the flip and source~3 is invisible
    \emph{before} it, and they carry restrictions even when the simplified answer is a polynomial with
    no denominator left to display them.
    \\[2pt]
    \textbf{Standards (2023 VA SOL):} \textbf{A2.EO.1a} --- add, subtract, multiply, or divide rational
    algebraic expressions, simplifying the result (today: multiply and divide). Applied throughout:
    \textbf{A2.EO.1b} (equivalent rational expressions with monomial and binomial factors; expressions
    limited to linear and quadratic) and \textbf{A2.EO.1d} (equivalence of forms). Builds on
    \textbf{A2.EO.3b} (factor completely, Lesson~3.2) and Lesson~4.1's restriction habit; it is the
    prerequisite for \textbf{A2.EO.1a/c} (4.3 --- a complex fraction \emph{is} a division) and for
    \textbf{A2.EI.4c} extraneous solutions (4.6).
\end{tcolorbox}

\renewcommand{\arraystretch}{1.4}
\begin{skillbox}[Priority Ideas \& Skills]{goldbox}
\begin{tabularx}{\linewidth}{@{}X|X@{}}
\textbf{Priority Ideas \& Skills} & \textbf{Key Understandings (the ``why'')} \\[2pt]
\hline
\vspace{-6pt}
\begin{itemize}[leftmargin=1.1em, nosep, after=\vspace{2pt}]
  \item \textbf{Factor everything first} --- the Unit~3 toolkit, now applied to four polynomials at once.
  \item \textbf{Multiply} by dividing out any numerator factor against any denominator factor, \emph{across the whole product}.
  \item \textbf{Divide} by \textbf{keep--change--flip} --- and flip \emph{before} cancelling.
  \item \textbf{List all three restriction sources}: both denominators \emph{and} the divisor's numerator.
  \item Handle \textbf{opposite binomials} ($\frac{a-b}{b-a}=-1$) inside a product.
  \item Connect back to 4.0/4.1: every restriction is a \textbf{hole} or a \textbf{wall} on the graph.
\end{itemize}
&
\vspace{-6pt}
\begin{itemize}[leftmargin=1.1em, nosep, after=\vspace{2pt}]
  \item An expanded product \emph{hides its own factors}. Multiplying out first destroys the only thing that makes the problem doable.
  \item Flipping is legal because $\frac cd\cdot\frac dc=1$ --- which is exactly why $c$ may not be zero.
  \item \textbf{Flipping trades one blind spot for another.} Read the restrictions off the problem as written, then add the divisor's numerator.
  \item You cannot divide by zero, and a fraction is zero exactly when its \emph{numerator} is.
  \item The $\div$ symbol is not a fraction bar --- nothing is lined up across it until you flip.
  \item An answer with no restriction list is half an answer, even when the answer is a polynomial.
\end{itemize}
\end{tabularx}
\end{skillbox}

\begin{skillbox}[{Vocabulary, Concepts, \& Theorems}]{sky}
\renewcommand{\arraystretch}{1.3}
\begin{tabularx}{\linewidth}{@{}l X@{}}
\textbf{Product rule} & \TallMath{\frac{a}{b}\cdot\frac{c}{d}=\frac{ac}{bd}} for $b\neq0$, $d\neq0$ --- but factor and divide out before you ever multiply. \\
\textbf{Reciprocal} & The fraction turned upside down: $\frac cd$ and $\frac dc$, whose product is $1$. It exists \emph{only} when $c\neq0$. \\
\textbf{Divisor} & The expression you are dividing \emph{by} --- the fraction after the $\div$ sign. \\
\textbf{Keep--change--flip} & \TallMath{\frac ab\div\frac cd=\frac ab\cdot\frac dc} --- keep the first, change the sign, flip the divisor. \\
\textbf{The three sources} & For $\frac AB\div\frac CD$: $B\neq0$, $D\neq0$, \emph{and} $C\neq0$. \\
\textbf{Hidden restriction} & An excluded value visible in neither the final answer nor (for source 3) the original denominators. \\
\textbf{Opposite binomials} & Same terms, every sign flipped; their quotient is $-1$ (Lesson 4.1, still in force). \\
\textbf{Hole vs.\ wall (4.0)} & Factor divides out $\Rightarrow$ removable discontinuity; factor stays $\Rightarrow$ vertical asymptote. \\
\end{tabularx}
\end{skillbox}

\begin{fixedskillbox}[Activate Prior Knowledge \& Spiral Review]{forestbg}
    The Warm-Up rehearses exactly three skills, one per leg of the lesson. \textbf{(1)} \emph{Fraction
    arithmetic on numbers}, where students already trust both operations: $\frac{6}{35}\cdot\frac{25}{9}$
    done the fast way (factor, divide out, \emph{then} multiply) establishes the discipline of the day,
    and $\frac49\div\frac{8}{15}$ names the \emph{reciprocal}. Item 1(c) plants today's genuinely new
    idea numerically before it is named: in $\frac49\div\frac{n}{15}$, the one value $n$ may not have is
    $0$ --- the divisor's \emph{numerator} is restricted. \textbf{(2)} \emph{Factoring completely}
    ($x^2-9$, $x^2+2x-8$, $x^2+3x$) --- all three reappear in the anchor. \textbf{(3)} \emph{Simplify
    and restrict} $\frac{x^2+3x}{x^2-9}$ with the hole/wall sort, carrying Lesson~4.1 forward intact.
    Debrief all three, but do \emph{not} explain where the divisor's restriction comes from --- the Hook
    does that.
\end{fixedskillbox}

\begin{skillbox}[Hook]{forestbg}
    Put two expressions on the board --- $Q(x)=\dfrac{x+1}{x-5}\div\dfrac{x-2}{x+3}$ and
    $R(x)=\dfrac{(x+1)(x+3)}{(x-5)(x-2)}$, with $R$ announced as ``$Q$ after keep--change--flip'' --- and
    have students fill a four-column table at $x=0,1,2,-3$. The first two columns agree
    ($\tfrac{3}{10}$ and $2$); at $x=2$ \emph{both} are undefined; at $x=-3$ they part company:
    $Q(-3)$ does not exist while $R(-3)=0$. Two questions run the discussion.
    \textbf{``Both die at $x=2$ --- for the same reason?''} (No: in $R$ the $(x-2)$ is a denominator; in
    $Q$ it sits on \emph{top} of the divisor and makes the whole divisor zero.) Then
    \textbf{``What broke at $x=-3$ in $Q$ that did not break in $R$?''} (The divisor itself is
    undefined --- and flipping quietly moved $(x+3)$ upstairs.) Land the sentence students will hear all
    period: \emph{flipping trades one blind spot for another}. Yesterday cancelling hid a restriction;
    today flipping hides a different one, so the restriction list is read off the problem \emph{as
    written} --- plus one more.
\end{skillbox}

\begin{skillbox}[Lesson]{forestbg}
\begin{multicols}{2}
\textbf{1. Multiplying: the rule is easy, the discipline is factoring.} State
$\frac ab\cdot\frac cd=\frac{ac}{bd}$ and then immediately tell students not to use it that way. The
Warm-Up already showed why on numbers; with polynomials the stakes are higher, because an
\emph{expanded} product hides its own factors and Unit~3 measured exactly what it costs to get them
back. Post the slogan: \textbf{factor everything first, divide out, multiply last --- and often not at
all.} Work the monomial case $\frac{6x^2}{5y}\cdot\frac{15y^3}{8x^5}=\frac{9y^2}{4x^3}$ ($x,y\neq0$),
noting that the restrictions came from \emph{both} denominators (A2.EO.1b's monomial clause).

\textbf{2. The four steps, one more denominator.} Same four steps as 4.1 --- factor, restrict, divide
out, write --- with step 2 now reading \emph{every} original denominator. The key new move in step 3 is
that cancelling goes \emph{across} the product: any numerator factor against any denominator factor.
Work the anchor
$\frac{x^2-9}{x^2+2x-8}\cdot\frac{x+4}{x^2+3x}=\frac{x-3}{x(x-2)}$, $x\neq0,2,-3,-4$, and stop on the
punch line: two of the four restrictions are now invisible, and they came from \emph{different}
fractions. Add the opposite-binomial case $\frac{x^2-25}{x+2}\cdot\frac{2x+4}{5-x}=-2(x+5)$
($x\neq-2,5$) --- yesterday's $-1$ has not gone anywhere.

\columnbreak

\textbf{3. Dividing --- and one warning worth the time.} Define the \textbf{reciprocal} and justify
keep--change--flip in one line: dividing by $\frac cd$ means multiplying by whatever undoes it, and
$\frac cd\cdot\frac dc=1$. Then make the students compute both cells of the wrong/right table:
$\frac43\div\frac54$ cancelled first gives $\frac{1}{15}$; flipped first gives $\frac{16}{15}$. The
$\div$ symbol is \emph{not} a fraction bar --- until you flip, nothing is lined up. Close with the
sentence that sets up Section~4: $\frac cd$ has a reciprocal only when $c\neq0$.

\textbf{4. The three sources (the lesson).} Build the table for $\frac AB\div\frac CD$: $B\neq0$ (the
first fraction must exist --- visible throughout), $D\neq0$ (the divisor must exist --- \emph{invisible
after} the flip, since $D$ moves upstairs), and $C\neq0$ (you cannot divide by zero, and the divisor is
zero exactly when its numerator is --- \emph{invisible before} the flip). Resolve the Hook with it, then
work the division anchor $\frac{x^2-4}{x^2+7x+12}\div\frac{x^2+2x}{x+3}=\frac{x-2}{x(x+4)}$,
$x\neq0,-2,-3,-4$, naming the source of each value aloud.

\textbf{5. The picture.} Finish on a graph, so today stays a \emph{function} lesson:
$P(x)=\frac{x^2-1}{x}\cdot\frac{x}{x+1}$ is the line $y=x-1$ with holes at $(0,-1)$ and $(-1,-2)$ ---
one hole per denominator. The simplified form has no denominator left to display them, and they are
excluded anyway. Name where this returns: \textbf{4.3} (a complex fraction is combine-then-divide),
\textbf{4.5} (every restriction becomes a hole or a wall), \textbf{4.6} (an extraneous solution is a
restriction sneaking back in).
\end{multicols}
\end{skillbox}

\begin{skillbox}[Explicit Instruction: Products and Quotients]{forestbg}
\begin{tabularx}{\linewidth}{@{}p{0.46\linewidth} X@{}}
\vspace{0pt}
\textbf{Steps (post these).}
\begin{enumerate}[leftmargin=1.3em, nosep, after=\vspace{2pt}]
  \item If it is a division, \textbf{keep--change--flip first}. Nothing cancels across a $\div$.
  \item Factor every numerator and denominator completely.
  \item \textbf{Write the restrictions now} --- both original denominators, \emph{plus} the divisor's
        numerator.
  \item Divide out any numerator factor against any denominator factor (watch for opposite
        binomials).
  \item Write the result, restrictions attached, left in factored form.
\end{enumerate}
&
\vspace{0pt}
\textbf{Worked, on the division anchor.}
\[ \frac{x^2-4}{x^2+7x+12}\div\frac{x^2+2x}{x+3}
   =\frac{(x-2)(x+2)}{(x+3)(x+4)}\cdot\frac{x+3}{x(x+2)} \]
Restrictions: $x\neq-3,-4$ (source 1), $x\neq-3$ (source 2), $x\neq0,-2$ (source 3).
\[ =\frac{x-2}{x(x+4)}, \qquad x\neq0,\;-2,\;-3,\;-4 \]
\textbf{The two questions students must answer aloud:} \emph{which fraction did each restriction come
from?} and \emph{which ones will your answer never remind you of?} A student who cannot answer the
second has not finished the problem.
\end{tabularx}
\end{skillbox}

\begin{skillbox}[Active Monitoring]{redbox}
Circulate during the guided practice and Tier work. Watch for and correct:
\begin{itemize}[leftmargin=1.2em, nosep]
  \item \textbf{multiplying out first} --- a quartic on the page means the problem just got harder; send them back to the factored line;
  \item \textbf{missing the divisor's-numerator restriction} --- the signature error of the lesson, and the one that returns in 4.3 and 4.6;
  \item \textbf{cancelling across the $\div$ before flipping} --- do not argue; hand the student $\frac x4\div\frac5x$ and the value $x=2$;
  \item \textbf{flipping the first fraction} instead of the divisor --- a mechanical slip one worked example fixes;
  \item \textbf{restrictions read off the final answer} --- yesterday's error, still costing points, and now worse because a product has two denominators;
  \item \textbf{no restrictions on a polynomial answer} --- $2x(x-4)$ still carries $x\neq0,-6$;
  \item \textbf{opposite binomials treated as identical}, losing the $-1$.
\end{itemize}
Cold-call prompts: ``Which \emph{factor} did you divide out, and which fraction was it in?'' ``Did you
flip before you cancelled --- prove it with the line above.'' ``Where did that restriction come from:
denominator one, denominator two, or the divisor's numerator?'' ``Your answer has no denominator ---
so why does it still have restrictions?''
\end{skillbox}

\begin{skillbox}[Group Work \& Differentiation]{redbox}
\textbf{Factor First, Flip Second, Restrict Always} activity (tiered; mirrors the notes).
\begin{multicols}{3}
\textbf{Tier R --- Remediate.} Five scaffolded items rising from a monomial product
($\frac{3x}{4}\cdot\frac{8}{9x^2}$, started for them) through a one-cancellation product
($\frac{x+2}{x-1}\cdot\frac{x-1}{x+5}$) and a factor-first product ($\frac{x^2-4}{x+3}\cdot\frac{2}{x-2}$)
to a first division ($\frac{x}{x+1}\div\frac{x-3}{x+1}$) whose restrictions are asked for by source.
Closes with the lesson's core question in plain words: $x=3$ was never a denominator --- why is it
excluded?

\columnbreak
\textbf{Tier A --- Approaching.} A four-row table (factored form / restrictions / simplest form) over
$\frac{10x^3}{3y^2}\cdot\frac{9y^5}{5x}$, $\frac{x^2+5x+6}{x^2-9}\cdot\frac{x-3}{x+2}$ (everything
divides out --- the answer is $\mathbf{1}$, and it still carries three restrictions),
$\frac{x^2-2x-8}{x^2-16}\div\frac{x+2}{x^2+4x}$ (answer $x$, carrying \emph{four} restrictions, two of
them from the divisor), and $\frac{9-x^2}{x+4}\cdot\frac{x^2+4x}{x-3}$ (\emph{opposite binomials}).
Then an \textbf{error analysis}: a student cancels the $x$'s across $\frac x4\div\frac5x$; groups
disprove it at $x=2$ ($\frac15$ vs.\ $\frac1{20}$), do it correctly, and state the rule in one
sentence.

\columnbreak
\textbf{Tier E --- Extension.} \emph{Part 1: restriction detective} --- simplify
$\frac{x^2-1}{x^2+5x+6}\div\frac{x^2-x}{x+3}$ to $\frac{x+1}{x(x+2)}$, then fill a table attributing
each of $x\neq0,1,-2,-3$ to its source(s) and marking whether it survives into the answer; one value
($-3$) has two sources. \emph{Part 2:} build a division whose answer is $\frac{x}{x-4}$ with
restrictions $x\neq4,-1,2$ --- the only way to get the $2$ is through the divisor's numerator.
\emph{Part 3:} recover a rectangle's length from its area and width by division, verify by multiplying
back, interpret the source-3 restriction geometrically (a rectangle with zero width), and find the
feasible domain as a union.
\end{multicols}
\end{skillbox}

\begin{skillbox}[Individual Work \& Assessment]{redbox}
\textbf{Exit Ticket} (independent, no notes): \textbf{(1)} simplify
$\frac{x^2-9}{x^2+2x}\div\frac{x-3}{x+2}$ showing the flipped, factored line; give
$\frac{x+3}{x}$ with \emph{all} restrictions $x\neq0,-2,3$, then name the two that are invisible in the
answer and where each came from; \textbf{(2)} \textbf{explain}, against a classmate's claim that
``once you flip, $x+2$ is on top, so $x=-2$ isn't a restriction anymore'' --- full credit requires both
halves: the restrictions are read from the problem as written, and flipping \emph{hides} a restriction
rather than removing it; \textbf{(3)} an \textbf{SOL-style multiple-choice} item on
$\frac{x^2-25}{x^2+3x}\div\frac{x-5}{x+3}$ whose distractors each break exactly one idea (restrictions
read off the answer / divisor's numerator forgotten / the wrong fraction flipped). Collect and sort
into ``got it / missing source 3 / restrictions from the answer / flipped the wrong fraction.'' The
source-3 pile decides whether 4.3 opens with a re-teach --- a complex fraction is a division, so the
same restriction reappears there immediately.
\end{skillbox}

\begin{skillbox}[Reinforcement \& Extension]{goldbox}
\textbf{Homework:} three monomial problems (one product, two quotients) including a case whose written
denominators are the constants $2$ and $5$, so that \emph{every} restriction comes from the divisor's
numerator; six binomial/quadratic problems spanning difference of squares, GCF-then-factor, a
polynomial answer that still carries restrictions ($2x(x-6)$), an \emph{opposite-binomial} product
($\frac{x^2-49}{x+1}\cdot\frac{3x+3}{7-x}=-3(x+7)$), and a quotient whose divisor's numerator is itself
a quadratic contributing two values; a rectangle model whose area simplifies to $\frac{x^2}{2}$ yet
keeps $x\neq-5$, with a feasible-domain union; and a justification item on
$\frac{x+1}{x-2}\div\frac{x+1}{x-3}$ versus $\frac{x-3}{x-2}$ --- equivalent, but with two extra holes
at $(3,0)$ and $\left(-1,\frac43\right)$. \textbf{Extension:} an error gallery --- three students,
three answers to $\frac{x^2-1}{x+4}\div\frac{x-1}{x^2-16}$ (correct / correct form with no restrictions
/ never flipped) --- and a \emph{design} task: build a division whose answer is $x+2$ with restrictions
$x\neq3,0,-1$, then describe the resulting line with three holes (a direct preview of 4.5).
\textbf{Preview:} Lesson~4.3, \emph{Adding \& Subtracting Rational Expressions} --- where a common
denominator finally becomes necessary, and where a \textbf{complex fraction} turns out to be nothing
but combine-then-divide, i.e.\ today's keep--change--flip with all three restriction sources intact.
\end{skillbox}
\end{document}
